%
% Main document
% ===========================================================================
% This is part of the document "Project documentation template".
% Authors: brd3, kaa1
%

%---------------------------------------------------------------------------
\documentclass[
	a4paper,					% paper format
	10pt,							% fontsize
	twoside,					% double-sided
	openright,				% begin new chapter on right side
	notitlepage,			% use no standard title page
	parskip=half,			% set paragraph skip to half of a line
]{scrreprt}					% KOMA-script report
%---------------------------------------------------------------------------

\raggedbottom
\KOMAoptions{cleardoublepage=plain}			% Add header and footer on blank pages


% Load Standard Packages:
%---------------------------------------------------------------------------
\usepackage[standard-baselineskips]{cmbright}

\usepackage[ngerman]{babel}										% german hyphenation
%\usepackage[latin1]{inputenc}  							% Unix/Linux - load extended character set (ISO 8859-1)
\usepackage[ansinew]{inputenc}  							% Windows - load extended character set (ISO 8859-1)
\usepackage[T1]{fontenc}											% hyphenation of words with �,� and �
\usepackage{textcomp}													% additional symbols
\usepackage{ae}																% better resolution of Type1-Fonts 
\usepackage{fancyhdr}													% simple manipulation of header and footer 
\usepackage{etoolbox}													% color manipulation of header and footer
\usepackage{graphicx}                      		% integration of images
\usepackage{float}														% floating objects
\usepackage{caption}													% for captions of figures and tables
\usepackage{booktabs}													% package for nicer tables
\usepackage{tocvsec2}													% provides means of controlling the sectional numbering
%---------------------------------------------------------------------------

% Load Math Packages
%---------------------------------------------------------------------------
\usepackage{amsmath}                    	   	% various features to facilitate writing math formulas
\usepackage{amsthm}                       	 	% enhanced version of latex's newtheorem
\usepackage{amsfonts}                      		% set of miscellaneous TeX fonts that augment the standard CM
\usepackage{amssymb}													% mathematical special characters
\usepackage{exscale}													% mathematical size corresponds to textsize
%---------------------------------------------------------------------------

% Package to facilitate placement of boxes at absolute positions
%---------------------------------------------------------------------------
\usepackage[absolute]{textpos}
\setlength{\TPHorizModule}{1mm}
\setlength{\TPVertModule}{1mm}
%---------------------------------------------------------------------------					
			
% Definition of Colors
%---------------------------------------------------------------------------
\RequirePackage{color}                          % Color (not xcolor!)
\definecolor{linkblue}{rgb}{0,0,0.8}            % Standard
\definecolor{darkblue}{rgb}{0,0.08,0.45}        % Dark blue
\definecolor{bfhgrey}{rgb}{0.41,0.49,0.57}      % BFH grey
%\definecolor{linkcolor}{rgb}{0,0,0.8}     			% Blue for the web- and cd-version!
\definecolor{linkcolor}{rgb}{0,0,0}        			% Black for the print-version!
%---------------------------------------------------------------------------

% Hyperref Package (Create links in a pdf)
%---------------------------------------------------------------------------
\usepackage[
	pdftex,ngerman,bookmarks,plainpages=false,pdfpagelabels,
	backref = {false},										% No index backreference
	colorlinks = {true},                  % Color links in a PDF
	hypertexnames = {true},               % no failures "same page(i)"
	bookmarksopen = {true},               % opens the bar on the left side
	bookmarksopenlevel = {0},             % depth of opened bookmarks
	pdftitle = {Template f�r Bachelor Thesis},	   	% PDF-property
	pdfauthor = {brd3},        					  % PDF-property
	pdfsubject = {LaTeX Template},        % PDF-property
	linkcolor = {linkcolor},              % Color of Links
	citecolor = {linkcolor},              % Color of Cite-Links
	urlcolor = {linkcolor},               % Color of URLs
]{hyperref}
%---------------------------------------------------------------------------
% Set up page dimension
%---------------------------------------------------------------------------
\usepackage{geometry}
\geometry{
	a4paper,
	left=28mm,
	right=15mm,
	top=30mm,
	headheight=20mm,
	headsep=10mm,
	textheight=242mm,
	footskip=15mm
}
%---------------------------------------------------------------------------

% Makeindex Package
%---------------------------------------------------------------------------
\usepackage{makeidx}                         		% To produce index
\makeindex                                    	% Index-Initialisation
%---------------------------------------------------------------------------

% Glossary Package
%---------------------------------------------------------------------------
% the glossaries package uses makeindex
% if you use TeXnicCenter do the following steps:
%  - Goto "Ausgabeprofile definieren" (ctrl + F7)
%  - Select the profile "LaTeX => PDF"
%  - Add in register "Nachbearbeitung" a new "Postprozessoren" point named Glossar
%  - Select makeindex.exe in the field "Anwendung" ( ..\MiKTeX x.x\miktex\bin\makeindex.exe )
%  - Add this [ -s "%tm.ist" -t "%tm.glg" -o "%tm.gls" "%tm.glo" ] in the field "Argumente"
%
% for futher informations go to http://ewus.de/tipp-1029.html
%---------------------------------------------------------------------------
\usepackage[nonumberlist]{glossaries}
\makeglossaries
\usepackage{xparse}
\DeclareDocumentCommand{\newdualentry}{ O{} O{} m m m m } {
	\newglossaryentry{gls-#3}{name={#5},text={#5\glsadd{#3}},
		description={#6},#1
	}
	\newacronym[type=\acronymtype, see={[Glossary:]{gls-#3}},#2]{#3}{#4}{#5\glsadd{gls-#3}}
}


% GLOSSARY
\newglossaryentry{imvr}{name={IMVR},description={Kurz f�r ``Images \& Music in VR''. Die Applikation, die im Rahmen dieses Projekts mit Unity 5 entwickelt wird}}
\newglossaryentry{unity}{name={Unity},description={Unity ist eine Spielengine, die das einfache Entwickeln von 3D-Applikation f�r diverse Endger�te erlaubt}}
\newglossaryentry{rift}{name={Oculus Rift},description={Ein \gls{hmd} von Oculus VR}}
\newglossaryentry{leap}{name={Leap Motion},description={Ein auf Infrarotkameras basierter Handerkennungs-Sensor}}


% ACRONYMS
\newacronym[type=\acronymtype]{hmd}{HMD}{Head-Mounted Display}
\newacronym{vr}{VR}{Virtual Reality}
%\newglossaryentry{ipd}{name={IPD},description={Beschreibt den Augenabstand.}}
\newdualentry{ipd}{IPD}{Inter-Pupillary Distance}{Beschreibt den Augenabstand und stellt ein wichtiges Mass f�r die stereoskopische Darstellung von Bildern dar}
%---------------------------------------------------------------------------

% Intro:
%---------------------------------------------------------------------------
\begin{document}                              	% Start Document
\settocdepth{section}														% Set depth of toc
\pagenumbering{roman}														
%---------------------------------------------------------------------------

\providecommand{\titel}{Bringing Your Hands Into Virtual Reality}		%  Hier den Titel des Berichts/Thesis eingeben					% Titel der Arbeit aus Datei titel.tex lesen
\providecommand{\versionnumber}{1.0}			%  Hier die aktuelle Versionsnummer eingeben
\providecommand{\versiondate}{11.06.2015}		%  Hier das Datum der aktuellen Version eingeben				% Versionsnummer und -datum aus Datei version.tex lesen

% Set up header and footer
%---------------------------------------------------------------------------
\makeatletter
\patchcmd{\@fancyhead}{\rlap}{\color{bfhgrey}\rlap}{}{}		% new color of header
\patchcmd{\@fancyfoot}{\rlap}{\color{bfhgrey}\rlap}{}{}		% new color of footer
\makeatother

\fancyhf{}																		% clean all fields
\fancypagestyle{plain}{												% new definition of plain style	
	\fancyfoot[OR,EL]{\footnotesize \thepage} 	% footer right part --> page number
	\fancyfoot[OL,ER]{\footnotesize \titel, Version \versionnumber, \versiondate}	% footer even page left part 
}

\renewcommand{\chaptermark}[1]{\markboth{\thechapter.  #1}{}}
\renewcommand{\headrulewidth}{0pt}				% no header stripline
\renewcommand{\footrulewidth}{0pt} 				% no bottom stripline

\pagestyle{plain}
%---------------------------------------------------------------------------


% Title Page and Abstract
%---------------------------------------------------------------------------
%%
% Project documentation template
% ===========================================================================
% This is part of the document "Project documentation template".
% Authors: brd3, kaa1
%

\begin{titlepage}


% BFH-Logo absolute placed at (28,12) on A4 
% Actually not a realy satisfactory solution but working.
%---------------------------------------------------------------------------
\setlength{\unitlength}{1mm}
\begin{textblock}{20}[0,0](28,12)
	\includegraphics[scale=1.0]{bilder/BFH_Logo_B.png}
\end{textblock}
\color{black}

% Institution / Titel / Untertitel / Autoren / Experten:
%---------------------------------------------------------------------------
\begin{flushleft}

\vspace*{21mm}

\fontsize{26pt}{40pt}\selectfont 
\titel 				\\							% Titel aus der Datei vorspann/titel.tex lesen
\vspace{2mm}

\fontsize{16pt}{24pt}\selectfont\vspace{0.3em}
Hier steht ein Untertitel 			\\							% Untertitel eingeben
\vspace{5mm}

\fontsize{10pt}{12pt}\selectfont
\textbf{Art der Arbeit (Semesterarbeit / Bachelorthesis / etc.)} \\									% eingeben
\vspace{7mm}

% Abstract (eingeben):
%---------------------------------------------------------------------------
\begin{textblock}{150}(28,100)
\fontsize{10pt}{12pt}\selectfont
[Kurztext (Abstract) einf�gen, falls gew�nscht] \\ 
Dieses Dokument dient als Vorlage f�r die Erstellung von Berichten nach den Richtlinien der BFH. Die Vorlage ist in \LaTeX{} erstellt und unterst�tzt das automatische Erstellen von diversen Verzeichnissen, Literaturangaben, Indexierung und Glossaren. Dieser kleine Text ist eine Zusammenfassung �ber das vorliegenden Dokument mit einer L�nge von 4 bis max. 8 Zeilen. \\
Das Titelbild kann in den Zeilen 157/158 der Datei template.tex ein- oder ausgeschaltet werden.
\end{textblock}

\begin{textblock}{150}(28,225)
\fontsize{10pt}{17pt}\selectfont
\begin{tabbing}
xxxxxxxxxxxxxxx\=xxxxxxxxxxxxxxxxxxxxxxxxxxxxxxxxxxxxxxxxxxxxxxx \kill
Studiengang:	\> [z.B. Elektro- und Kommunikationstechnik]	\\			% Namen eingeben
Autoren:		\> [Test Peter, M�ster R�s�]		\\					% Namen eingeben
Betreuer:	\> [Dr.~Xxxx Xxxx, Dr.~Yyyy Yyyy]		\\					% Namen eingeben
Auftraggeber:	\> [Wwwww AG]						\\					% Namen eingeben
Experten:		\> [Dr.~Zzzz Zzzz]				\\					% Namen eingeben
Datum:			\> \versiondate					\\		% aus Datei vorspann/version.tex lesen
\end{tabbing}

\end{textblock}
\end{flushleft}

\begin{textblock}{150}(28,280)
\noindent 
\color{bfhgrey}\fontsize{9pt}{10pt}\selectfont
Berner Fachhochschule | Haute �cole sp�cialis�e bernoise | Bern University of Applied Sciences
\color{black}\selectfont
\end{textblock}


\end{titlepage}

%
% ===========================================================================
% EOF
%
		% activate for Titelseite ohne Bild
%
% Project documentation template
% ===========================================================================
% This is part of the document "Project documentation template".
% Authors: brd3, kaa1
%

\begin{titlepage}


% BFH-Logo absolute placed at (28,12) on A4 and picture (16:9 or 15cm x 8.5cm)
% Actually not a realy satisfactory solution but working.
%---------------------------------------------------------------------------
\setlength{\unitlength}{1mm}
\begin{textblock}{20}[0,0](28,12)
	\includegraphics[scale=1.0]{bilder/BFH_Logo_B.png}
\end{textblock}

\begin{textblock}{154}(28,48)
	\begin{picture}(150,2)
		\put(0,0){\color{bfhgrey}\rule{150mm}{2mm}}
	\end{picture}
\end{textblock}

\begin{textblock}{154}[0,0](28,55)
	\centering
	\includegraphics[width=0.7\linewidth]{bilder/logo}			% Titelbild definieren
\end{textblock}

\begin{textblock}{154}(28,155)
	\begin{picture}(150,2)
		\put(0,0){\color{bfhgrey}\rule{150mm}{2mm}}
	\end{picture}
\end{textblock}
\color{black}

% Institution / Titel / Untertitel / Autoren / Experten:
%---------------------------------------------------------------------------
\begin{flushleft}

\vspace*{125mm}

\fontsize{26pt}{28pt}\selectfont 
\titel 				\\							% Titel aus der Datei vorspann/titel.tex lesen
\vspace{2mm}

\fontsize{16pt}{20pt}\selectfont\vspace{0.3em}
Schlussbericht 			\\							% Untertitel eingeben
\vspace{5mm}

\fontsize{10pt}{12pt}\selectfont
%\textbf{Art der Arbeit (Semesterarbeit / Bachelorthesis / etc.)} \\									% eingeben
\vspace{3mm}

% Abstract (eingeben):
%---------------------------------------------------------------------------
%\begin{textblock}{150}(28,190)
%\fontsize{10pt}{12pt}\selectfont
%[Kurztext (Abstract) einf�gen, falls gew�nscht] \\ 
%Dieses Dokument dient als Vorlage f�r die Erstellung von Berichten nach den Richtlinien der BFH. Die Vorlage ist in \LaTeX{} erstellt und unterst�tzt das automatische Erstellen von diversen Verzeichnissen, Literaturangaben, Indexierung und Glossaren. Dieser kleine Text ist eine Zusammenfassung �ber das vorliegenden Dokument mit einer L�nge von 4 bis max. 8 Zeilen. \\
%Das Titelbild kann in den Zeilen 157/158 der Datei template.tex ein- oder ausgeschaltet werden.
%\end{textblock}

\begin{textblock}{150}(28,225)
\fontsize{10pt}{17pt}\selectfont
\begin{tabbing}
xxxxxxxxxxxxxxx\=xxxxxxxxxxxxxxxxxxxxxxxxxxxxxxxxxxxxxxxxxxxxxxx \kill
Studiengang:	\> Informatik	\\			% Namen eingeben
Autoren:		\> Simon Meer		\\					% Namen eingeben
Betreuer:	\> Prof. Urs K�nzler		\\					% Namen eingeben
%Auftraggeber:	\> [Wwwww AG]						\\					% Namen eingeben
Experten:		\> Yves Petitpierre				\\					% Namen eingeben
Datum:			\> \versiondate					\\		% aus Datei vorspann/version.tex lesen
\end{tabbing}

\end{textblock}
\end{flushleft}

\begin{textblock}{150}(28,280)
\noindent 
\color{bfhgrey}\fontsize{9pt}{10pt}\selectfont
Berner Fachhochschule | Haute �cole sp�cialis�e bernoise | Bern University of Applied Sciences
\color{black}\selectfont
\end{textblock}


\end{titlepage}

%
% ===========================================================================
% EOF
%
			% activate for Titelseite mit Bild
% Versionenkontrolle :
% -----------------------------------------------
%\pagebreak
\begin{textblock}{180}(15,150)
\color{black}
\begin{huge}
Versionen
\end{huge}
\vspace{10mm}

\fontsize{10pt}{18pt}\selectfont
\begin{tabbing}
xxxxxxxxxxx\=xxxxxxxxxxxxxxx\=xxxxxxxxxxxxxx\=xxxxxxxxxxxxxxxxxxxxxxxxxxxxxxxxxxxxxxxxxxxxxxx \kill
Version	\> Datum	\> Status		\> Bemerkungen		\\
0.9	\> 14.04.2015	\> Entwurf		\> Outline erstellt	\\
\end{tabbing}

\end{textblock}
%\pagebreak
\cleardoubleemptypage
\setcounter{page}{1}
%\cleardoublepage
%\phantomsection 
%\addcontentsline{toc}{chapter}{Management Summary}
%\chapter*{Management Summary}
\label{chap:managementSummary}

Lorem ipsum dolor sit amet, consectetur adipiscing elit. Phasellus scelerisque, leo sed iaculis ornare, mi leo semper urna, ac elementum libero est at risus. Donec eget aliquam urna. Lorem ipsum dolor sit amet, consectetur adipiscing elit. Nunc fermentum nunc sollicitudin leo porttitor volutpat. Duis ac enim lectus, quis malesuada lectus. Aenean vestibulum suscipit justo, in suscipit augue venenatis a. Donec interdum nibh ligula. Aliquam vitae dui a odio cursus interdum quis vitae mi. Phasellus ornare tortor fringilla velit accumsan quis tincidunt magna eleifend. Praesent nisl nibh, cursus in mattis ac, ultrices ac nulla. Nulla ante urna, aliquet eu tempus ut, feugiat id nisl. Nunc sit amet mauris vitae turpis scelerisque mattis et sed metus. Aliquam interdum congue odio, sed semper elit ullamcorper vitae. Morbi orci elit, feugiat vel hendrerit nec, sollicitudin non massa. Quisque lacus metus, vulputate id ullamcorper id, consequat eget orci \nocite{kopka:band1} \nocite{Marti06}. 

\cleardoubleemptypage
%---------------------------------------------------------------------------

% Table of contents
%---------------------------------------------------------------------------
\tableofcontents
\cleardoublepage
%---------------------------------------------------------------------------

% Main part:
%---------------------------------------------------------------------------
\pagenumbering{arabic}

\chapter{Einleitung}
\label{chap:einleitung}

Der Markt der stereoskopischen Brillen befindet sich in einem starken Wachstum. Oculus Rift, SteamVR, Project Morpheus, Nvidia VR - die Liste der teilnehmenden Partien w�chst und w�chst, doch neue Ausgabeger�te ben�tigen auch neue Eingabeger�te. Die Leap Motion kann als ein solches gez�hlt werden und bietet eine freie Handerkennung ohne irgendwelche Handschuhe tragen zu m�ssen.

In diesem Projekt geht es darum, die neuartige Eingabem�glichkeiten der Leap Motion in Verbindung mit der Oculus Rift einzusetzen und damit eine interaktive \gls{vr}-Applikation zu schreiben. Im Rahmen eines vorhergehenden Projekts wurde die Materie bereits bearbeitet und die M�glichkeiten erforscht. Es geht nun also haupts�chlich um die erfolgreiche Anwendung der erarbeiteten Informationen.
\include{kapitel/projektstruktur}
\include{kapitel/zielbestimmung}
\chapter{Schlussfolgerungen/Fazit}
\label{chap:schlussfolgerungen}

Tr�stlicheres. Tod Baus treuh�nderischem Feldern ade war Wetten peu exaltiert, sondern just wiegen Welchen seh Egels all Eis Namur wutentbranntes Aas. Creme hob Herklit abnimmt ruh eben Auto Tr�ffels dept Storchs tr�gten Axt Alarmen Malereien Acker gar Puder Bea wohlmeinendere ba� Autovermietung Pistole, Bern edel tanze Ergusses bin kam defektes, Gag wedle Franziska Ehe ja unsriger hob gabt. Fells �berregionalen Eklat droht Dr spe Barden Boy gib Frl Sonette. Tito fesseln sich ade Big eng Julis lobe Gas auf F�rberei folgen Extension Brandmal stillte C. Wartens half Box umgehauter umworbenes Bruchst�cken, tov Ehe Pokals geh tapsige, segnete sag Eink�ufe wer Aas weh einzahlendes H�geln. Heft abschn�rend Bandit dm dies l�gen tankte hat.Abeter teilt geize Bzw turne mystisch G�thes Dorfes, Cha Beo Deuterium Allergien Bar von gekrochen ahndest. Art falls gehe Gef�ss ortest fair, ade adlige klarste was tolle Ada Obmann, gen C. Klausel nage allm�chtig. Zweitem. Sattle Gebote Droht B�e F�chers.

\newpage

Welchen seh Egels all Eis Namur wutentbranntes Aas. Creme hob Herklit abnimmt ruh eben Auto Tr�ffels dept Storchs tr�gten Axt Alarmen Malereien Acker gar Puder Bea wohlmeinendere ba� Autovermietung Pistole, Bern edel tanze Ergusses bin kam defektes, Gag wedle Franziska Ehe ja unsriger hob gabt. Fells �berregionalen Eklat droht Dr spe Barden Boy gib Frl Sonette. Tito fesseln sich ade Big eng Julis lobe Gas auf F�rberei folgen Extension Brandmal stillte C. Wartens half Box umgehauter umworbenes Bruchst�cken, tov Ehe Pokals geh tapsige, segnete sag Eink�ufe wer Aas weh einzahlendes H�geln. Heft abschn�rend Bandit dm dies l�gen tankte hat.Abeter teilt geize Bzw turne mystisch G�thes Dorfes, Cha Beo Deuterium Allergien Bar von gekrochen ahndest. Art falls gehe Gef�ss ortest fair, ade adlige klarste was tolle Ada Obmann, gen C. Klausel nage allm�chtig. 

%---------------------------------------------------------------------------

% Selbst�ndigkeitserkl�rung
%---------------------------------------------------------------------------
%\cleardoublepage
%\phantomsection 
%\addcontentsline{toc}{chapter}{Selbst�ndigkeitserkl�rung}
%\chapter*{Selbst�ndigkeitserkl�rung}
\label{chap:selbstaendigkeitserklaerung}

\vspace*{10mm} 

Ich/wir best�tige/n, dass ich/wir die vorliegende Arbeit selbstst�ndig und ohne Benutzung anderer als der im Literaturverzeichnis angegebenen Quellen und Hilfsmittel angefertigt habe/n. S�mtliche Textstellen, die nicht von mir/uns stammen, sind als Zitate gekennzeichnet und mit dem genauen Hinweis auf ihre Herkunft versehen. 

\vspace{15mm}

\begin{tabbing}
xxxxxxxxxxxxxxxxxxxxxxxxx\=xxxxxxxxxxxxxxxxxxxxxxxxxxxxxx\=xxxxxxxxxxxxxxxxxxxxxxxxxxxxxx\kill
Ort, Datum:		\> [Biel/Burgdorf], \versiondate \\ \\ 
Namen Vornamen:	\> [Test Peter] 	\> [M�ster R�s�] \\ \\ \\ \\ 
Unterschriften:	\> ......................................\> ...................................... \\
\end{tabbing}

%---------------------------------------------------------------------------

% Glossary
%---------------------------------------------------------------------------
\cleardoublepage
\phantomsection 
\addcontentsline{toc}{chapter}{Glossar}
\renewcommand{\glossaryname}{Glossar}
\printglossary
%---------------------------------------------------------------------------

% Bibliography
%---------------------------------------------------------------------------
\cleardoublepage
\phantomsection 
\addcontentsline{toc}{chapter}{Literaturverzeichnis}
\bibliographystyle{IEEEtranS}
\bibliography{datenbanken/bibliography}{}
%---------------------------------------------------------------------------

% Listings
%---------------------------------------------------------------------------
%\cleardoublepage
%\phantomsection 
%\addcontentsline{toc}{chapter}{Abbildungsverzeichnis}
%\listoffigures
%\cleardoublepage
%\phantomsection 
%\addcontentsline{toc}{chapter}{Tabellenverzeichnis}
%\listoftables
%---------------------------------------------------------------------------

% Index
%---------------------------------------------------------------------------
\cleardoublepage
\phantomsection 
\addcontentsline{toc}{chapter}{Stichwortverzeichnis}
\renewcommand{\indexname}{Stichwortverzeichnis}
\printindex
%---------------------------------------------------------------------------

% Attachment:
%---------------------------------------------------------------------------
\appendix
\settocdepth{section}
\chapter{Beliebiger Anhang}
\label{chap:bel_anhang}

Phasellus eget velit massa, sed faucibus nisi. Etiam tincidunt libero viverra lorem bibendum ut rutrum nisi volutpat. Donec non quam vitae lacus egestas suscipit at eu nisi. Maecenas non orci risus, at egestas tellus. Vivamus quis est pretium mauris fermentum consectetur. Cras non dolor vitae nulla molestie facilisis. Aliquam euismod nisl eget risus pretium non suscipit nulla feugiat. Nam in tortor sapien. Nam lectus nibh, laoreet eu ultrices nec, consequat nec sem. Nulla leo turpis, suscipit in vulputate a, dapibus molestie quam. Vestibulum pretium, purus sed suscipit tempus, turpis purus fermentum diam, id cursus enim mi a tortor. Proin imperdiet varius pellentesque. Nam congue, enim sit amet iaculis venenatis, dui neque ornare purus, laoreet porttitor nunc justo vel velit. Suspendisse potenti. Nulla facilisi.

\chapter{Weiterer Anhang}
\label{chap:anhang_B}

\section{Test 1}
Phasellus eget velit massa, sed faucibus nisi. Etiam tincidunt libero viverra lorem bibendum ut rutrum nisi volutpat. Donec non quam vitae lacus egestas suscipit at eu nisi. Maecenas non orci risus, at egestas tellus. Vivamus quis est pretium mauris fermentum consectetur. Cras non dolor vitae nulla molestie facilisis. Aliquam euismod nisl eget risus pretium non suscipit nulla feugiat. Nam in tortor sapien. 

\subsection{Umfeld}
Nam lectus nibh, laoreet eu ultrices nec, consequat nec sem. Nulla leo turpis, suscipit in vulputate a, dapibus molestie quam. Vestibulum pretium, purus sed suscipit tempus, turpis purus fermentum diam, id cursus enim mi a tortor. Proin imperdiet varius pellentesque. Nam congue, enim sit amet iaculis venenatis, dui neque ornare purus, laoreet porttitor nunc justo vel velit. Suspendisse potenti. Nulla facilisi.

\chapter{Inhalt der CD-ROM}
\label{chap:Inhalt_CDROM}

Inhaltsverzeichnis der beiliegenden CD-ROM, ev. Verzeichnisbaum, etc.
%---------------------------------------------------------------------------

%---------------------------------------------------------------------------
\end{document}

